\par\Needspace{6\baselineskip}\medskip
\hypertarget{native-209-article155}{}\pdfbookmark[2]{155.}{article-155}
\begin{center}155.\end{center}
Damit deutlicher hervortritt, wie genau die gefundenen Elemente sind, berechnen wir aus ihnen den mittleren Ort. Für Okt.\ $17{,}415\,011$ ergibt sich die mittlere Anomalie $=332^{0}\,28'\,54''{,}76$, daraus die wahre Anomalie $315^{0}\,1'\,23''{,}02$ und $\log r'=0{,}325\,9877$ (siehe \SourceLeaf{210}die Beispiele in den Artikeln 13, 14). Diese wahre Anomalie müsste gleich der wahren Anomalie bei der ersten Beobachtung, vermehrt um den Winkel $2f''$, oder der wahren Anomalie bei der dritten Beobachtung, vermindert um den Winkel $2f$, sein, d.\ h.\ $=315^{0}\,1'\,22''{,}98$. Der Logarithmus des Radiusvektors müsste dagegen $=0{,}325\,9878$ sein: Die Unterschiede sind zu vernachlässigen. Führt man die Rechnung für die mittlere Beobachtung bis zum geozentrischen Ort fort, so weichen die erhaltenen Zahlen von der Beobachtung nur um wenige Hundertstel einer Bogensekunde ab (Artikel 63). Solche Unterschiede gehen gewissermaßen in den unvermeidlichen Fehlern auf, die aus der begrenzten Genauigkeit der Tafeln entstehen.

Wir haben das vorhergehende Beispiel mit größter Genauigkeit durchgerechnet, um zu zeigen, wie leicht sich mit unserer Methode eine möglichst genaue Lösung gewinnen lässt. In der Praxis wird man diesem Beispiel nur selten so peinlich genau folgen müssen: Meist genügt es, durchweg \textit{sechs} Dezimalstellen zu verwenden. In unserem Beispiel hätte schon die zweite Hypothese eine ebenso große, die erste eine mehr als ausreichende Genauigkeit geliefert. Wir glauben, dass die Leser gerne die aus der dritten Hypothese ermittelten Elemente mit denen vergleichen werden, die sich bei gleicher Anwendung der zweiten oder gar der ersten Hypothese für denselben Zweck ergeben hätten. Diese drei Sätze von Elementen stellen wir in der folgenden Tabelle zusammen:
\begin{center}
\begingroup\footnotesize\setlength{\tabcolsep}{2pt}\renewcommand{\arraystretch}{1.2}
\begin{tabular}{|p{1.35in}|r|r|r|}\hline
& aus Hypothese III & aus Hypothese II & aus Hypothese I\\\hline
Epoche der mittl.\ Länge 1805 & $41^{0}\,52'\,21''{,}68$ & $41^{0}\,52'\,18''{,}40$ & $42^{0}\,12'\,37''{,}83$\\
mittlere tägl.\ Bewegung & $824''{,}7990$ & $824''{,}7978$ & $823''{,}5026$\\
Perihel & $52^{0}\,18'\,9''{,}30$ & $52^{0}\,18'\,6''{,}66$ & $52^{0}\,41'\,9''{,}81$\\
$\varphi$ & $14\quad12\quad1{,}87$ & $14\quad11\quad59{,}94$ & $14\quad24\quad27{,}49$\\
Log.\ der großen Halbachse & $0{,}422\,4389$ & $0{,}422\,4394$ & $0{,}422\,8944$\\
aufsteigender Knoten & $171^{0}\,7'\,48''{,}73$ & $171^{0}\,7'\,49''{,}15$ & $171^{0}\,5'\,48''{,}86$\\
Bahnneigung & $13\quad6\quad44{,}10$ & $13\quad6\quad45{,}12$ & $13\quad2\quad37{,}50$\\\hline
\end{tabular}\endgroup
\end{center}

Berechnet man den heliozentrischen Ort in der Bahn für die mittlere Beobachtung mit dem zweiten Satz von Elementen, so erhält man für den Logarithmus des Radiusvektors den Fehler $=0$ und für die Länge in der Bahn den Fehler $=0''{,}03$. Berechnet man diesen Ort hingegen mit dem aus der ersten Hypothese abgeleiteten Satz, so erhält man für den Logarithmus des Radiusvektors den Fehler $=0{,}000\,0002$ und für die Länge in der Bahn den Fehler $=1''{,}31$. Führt man die Rechnung bis zum geozentrischen Ort fort, so erhält man
\SourceLeaf{211}
\begin{center}
\begingroup\small\setlength{\tabcolsep}{4pt}\renewcommand{\arraystretch}{1.2}
\begin{tabular}{|l|r|r|}\hline
& aus Hypothese II & aus Hypothese I\\\hline
Geozentrische Länge & $352^{0}\,34'\,22''{,}26$ & $352^{0}\,34'\,19''{,}97$\\
Fehler & $0{,}14$ & $2{,}15$\\
geozentrische Breite & $-6\quad21\quad55{,}06$ & $-6\quad21\quad54{,}47$\\
Fehler & $0{,}01$ & $0{,}60$\\\hline
\end{tabular}\endgroup
\end{center}
